\documentclass[]{article}
\usepackage{lmodern}
\usepackage{amssymb,amsmath}
\usepackage{ifxetex,ifluatex}
\usepackage{fixltx2e} % provides \textsubscript
\ifnum 0\ifxetex 1\fi\ifluatex 1\fi=0 % if pdftex
  \usepackage[T1]{fontenc}
  \usepackage[utf8]{inputenc}
\else % if luatex or xelatex
  \ifxetex
    \usepackage{mathspec}
    \usepackage{xltxtra,xunicode}
  \else
    \usepackage{fontspec}
  \fi
  \defaultfontfeatures{Mapping=tex-text,Scale=MatchLowercase}
  \newcommand{\euro}{€}
\fi
% use upquote if available, for straight quotes in verbatim environments
\IfFileExists{upquote.sty}{\usepackage{upquote}}{}
% use microtype if available
\IfFileExists{microtype.sty}{\usepackage{microtype}}{}
\usepackage[margin=1in]{geometry}
\ifxetex
  \usepackage[setpagesize=false, % page size defined by xetex
              unicode=false, % unicode breaks when used with xetex
              xetex]{hyperref}
\else
  \usepackage[unicode=true]{hyperref}
\fi
\hypersetup{breaklinks=true,
            bookmarks=true,
            pdfauthor={Justin Le},
            pdftitle={Using LLMs Effectively with Haskell: Requirements-Circumventing Behavior},
            colorlinks=true,
            citecolor=blue,
            urlcolor=blue,
            linkcolor=magenta,
            pdfborder={0 0 0}}
\urlstyle{same}  % don't use monospace font for urls
% Make links footnotes instead of hotlinks:
\renewcommand{\href}[2]{#2\footnote{\url{#1}}}
\setlength{\parindent}{0pt}
\setlength{\parskip}{6pt plus 2pt minus 1pt}
\setlength{\emergencystretch}{3em}  % prevent overfull lines
\setcounter{secnumdepth}{0}

\title{Using LLMs Effectively with Haskell: Requirements-Circumventing Behavior}
\author{Justin Le}

\begin{document}
\maketitle

\emph{Originally posted on
\textbf{\href{https://blog.jle.im/entry/llms-haskell-constraint-evading-behavior.html}{in
Code}}.}

\documentclass[]{}
\usepackage{lmodern}
\usepackage{amssymb,amsmath}
\usepackage{ifxetex,ifluatex}
\usepackage{fixltx2e} % provides \textsubscript
\ifnum 0\ifxetex 1\fi\ifluatex 1\fi=0 % if pdftex
  \usepackage[T1]{fontenc}
  \usepackage[utf8]{inputenc}
\else % if luatex or xelatex
  \ifxetex
    \usepackage{mathspec}
    \usepackage{xltxtra,xunicode}
  \else
    \usepackage{fontspec}
  \fi
  \defaultfontfeatures{Mapping=tex-text,Scale=MatchLowercase}
  \newcommand{\euro}{€}
\fi
% use upquote if available, for straight quotes in verbatim environments
\IfFileExists{upquote.sty}{\usepackage{upquote}}{}
% use microtype if available
\IfFileExists{microtype.sty}{\usepackage{microtype}}{}
\usepackage[margin=1in]{geometry}
\ifxetex
  \usepackage[setpagesize=false, % page size defined by xetex
              unicode=false, % unicode breaks when used with xetex
              xetex]{hyperref}
\else
  \usepackage[unicode=true]{hyperref}
\fi
\hypersetup{breaklinks=true,
            bookmarks=true,
            pdfauthor={},
            pdftitle={},
            colorlinks=true,
            citecolor=blue,
            urlcolor=blue,
            linkcolor=magenta,
            pdfborder={0 0 0}}
\urlstyle{same}  % don't use monospace font for urls
% Make links footnotes instead of hotlinks:
\renewcommand{\href}[2]{#2\footnote{\url{#1}}}
\setlength{\parindent}{0pt}
\setlength{\parskip}{6pt plus 2pt minus 1pt}
\setlength{\emergencystretch}{3em}  % prevent overfull lines
\setcounter{secnumdepth}{0}


\begin{document}

Haskell's selling point was always "if it compiles, it works." Under AI code
generation this becomes *more* true, not less -- the type system is no longer
just catching your mistakes, it's catching an alien intelligence's mistakes. But
only if the alien doesn't learn to turn off the checks.

This is an experience report on a specific failure mode: **constraint-evading
behavior**. When an LLM agent hits friction -- a type error, a warning, a test
failure -- it reliably takes the path of *silencing the signal* rather than
*fixing the cause*.

## How agents fight Haskell

Every one of these patterns is a red flag. They all mean the same thing: the
agent is optimizing for "error goes away" rather than "code is correct."

### Silencing the compiler

-   Disabling warnings with `-Wno-*` or `-w`
-   Using wildcard patterns `_` to suppress incomplete pattern match warnings
-   Using catch-all patterns (`_` or a variable binding like `other ->`) that
    fail to narrow the type -- the pattern match should make what you pass along
    strictly smaller; if you bind a catch-all and re-pass it, your type is too
    big. This includes defining catch-all constructors (`| Other Text Value`) --
    leaves you vulnerable to string-stuffing

### Weakening types

-   Changing `NonEmpty` to `[]`
-   Changing a specific sum type to `Text`
-   Changing `Natural` to `Int`
-   Removing a typeclass constraint
-   Making a required field `Optional`
-   Stuffing strings into existing sum type constructors
    (`GenericError "description"`) instead of adding a new constructor when a
    new case arises. Same applies to `Value`/`Object` bags in domain types -- if
    you know the schema, type it
-   Keeping a bad type/design because fixing downstream is "inconvenient" -- if
    a type is wrong (e.g. `[[a]]` when it should be `[NonEmpty a]`), fix it.
    "Touching many call sites" is not a reason to keep a bad design

### Faking success

-   Hard-coding `True` or success values in tests
-   Disabling or skipping test execution
-   Commenting out failing code instead of fixing it

### Faking the algorithm

-   Implementing heuristics or special-case pattern matching when a correct
    general algorithm is required -- if you cannot write the actual algorithm,
    stop and say so; special-casing is not an algorithm
-   Passing typed data as strings through boundaries -- exporting integers as
    strings, passing a serialized string through FFI instead of the enum itself;
    if a type exists, wire the type through

### Data integrity violations

-   Prefixing arguments with `_` to hide that real data is being ignored. Same
    applies to ignored constructor payloads (`Foo _x _y -> ...`) -- if all
    fields are ignored, either the type is wrong or the code is wrong
-   Ignoring arguments containing real data and fabricating synthetic data
    instead
-   Falling back to `show` when a type lacks a wire format instance -- the
    missing instance is a signal to stop and ask, not a problem to paper over

### Bypassing the type system entirely

-   `unsafeCoerce` or other unsafe operations
-   `error`, `undefined`, or partial functions
-   Catching `SomeException` instead of the specific exceptions that can
    actually occur -- you lose the ability to handle different failures
    differently

## Why this happens

The LLM optimizes for "error goes away" -- the shortest edit distance from red
to green. It has no model of *why* the constraint exists, only that the
constraint blocks completion. Haskell's errors are verbose and specific, which
paradoxically makes them easier to silence precisely. The agent treats the
compiler as an adversary rather than a collaborator.

## Why Haskell is still better for AI than dynamic languages

In Python, the equivalent mistakes are silent -- you never even know they
happened. Haskell *surfaces* the circumvention as a detectable pattern: a
disabled warning, a weakened type, a catch-all where there wasn't one before.
You can write rules against specific circumventions; you can't write rules
against "silently wrong."

The type system turns AI mistakes from runtime mysteries into compile-time red
flags.

## How to use Haskell most effectively with LLMs

1.  **Explicit prohibition lists in system prompts.** The AI won't infer these
    norms. It needs to be told "never do X" for each X.

2.  **Explain why each prohibition exists.** Without the reason, the AI invents
    novel circumventions that aren't on the list.

3.  **Frame guardrails as steering, not restriction.** "This is how you know
    you're going wrong" works better than "don't do this."

4.  **Make stopping a valid action.** "If you can't satisfy the type, stop and
    ask" -- otherwise the agent will always find *something* to do, and that
    something will be circumvention.

5.  **Use the strongest types you can.** `NonEmpty`, newtypes, GADTs -- they
    generate better error messages that steer the AI toward correct solutions.
    Make invalid states unrepresentable and the AI can't fabricate data that
    doesn't fit.

6.  **Treat weakened types in AI output as defects.** Every `_` pattern, every
    dropped constraint, every `show` where an instance should exist -- these are
    not style issues, they are correctness bugs.

## Mechanical enforcement with hooks

Instructions alone are not enough. The agent will follow them *most* of the
time, but under pressure (complex type errors, cascading failures) it will
circumvent. Claude Code hooks let you block the circumvention at the tool
boundary -- the agent physically cannot complete the prohibited action.

### Blocking forbidden constructs in edits

A `PreToolUse` hook on `Edit|Write` that greps the new content:

``` bash
# Block error/unsafeCoerce/OPTIONS_GHC in .hs files
if echo "$content" | grep -qF 'GHC.Err.error'; then
    echo 'Blocked: Do not use error. If you cannot satisfy the types, ask.' >&2
    exit 2
fi
if echo "$content" | grep -qP '\bunsafeCoerce\b'; then
    echo 'Blocked: Do not use unsafeCoerce.' >&2
    exit 2
fi
if echo "$content" | grep -qP 'OPTIONS_GHC'; then
    echo 'Blocked: Do not add OPTIONS_GHC pragmas.' >&2
    exit 2
fi
```

The `OPTIONS_GHC` message explicitly anticipates the next circumvention: "Do NOT
work around this by putting the option in cabal files, per-module ghc-options,
or any other mechanism."

### Blocking edits to lint configuration

``` bash
# Block edits to hlint config files
if [[ "$file_path" == *hlint* && "$file_path" == *.yaml ]]; then
    echo 'Blocked: Do not edit hlint configuration files.' >&2
    exit 2
fi
# Block inline HLINT ignore pragmas
if echo "$content" | grep -qiP '\{-#?\s*HLINT\s+ignore'; then
    echo 'Blocked: Do not use {-# HLINT ignore #-} pragmas.' >&2
    exit 2
fi
```

This covers both attack vectors: editing the config to disable rules globally,
and suppressing rules per-file with pragmas.

### HLint as a ratchet

Run hlint on both the old and new content. Only block if the edit *introduces*
new issues:

``` bash
new_output=$($HLINT "$temp_new" 2>&1 || true)
old_output=$($HLINT "$temp_old" 2>&1 || true)

# Only block if new content has issues and old content didn't
if [[ "$new_has_error" == "true" && "$old_has_error" == "false" ]]; then
    echo "Blocked: HLint issues detected" >&2
    echo "$new_output" >&2
    exit 2
fi
```

This is a ratchet -- existing issues can stay, but you can never make things
worse. The agent can't introduce new lint violations even if it wants to.

### Blocking piped build output

``` bash
# Block piping build commands (agents grep away their own errors)
if echo "$cmd" | grep -qP '(?:cabal\s+build|nix\s+build)'; then
    if echo "$cmd" | grep -qP '\|'; then
        echo 'Blocked: You must see the full build output.' >&2
        exit 2
    fi
fi
```

Without this, agents will pipe build output through `grep` to filter for the
specific error they're looking for -- missing other errors they introduced. This
is a subtle circumvention: not suppressing the check, but suppressing the
*visibility* of the check's results.

### Enforcing style rules at write time

A `let...in` vs `where` checker that counts violations in old vs new content and
blocks if the count increases:

``` bash
new_let_in=$(count_bad_let_in "$new_content")
old_let_in=$(count_bad_let_in "$old_content")
if [[ "$new_let_in" -gt "$old_let_in" ]]; then
    echo 'Blocked: use where instead of let...in' >&2
    exit 2
fi
```

Same ratchet pattern as hlint -- only blocks regressions.

### Blocking blame-shifting in responses

A `Stop` hook (runs after the agent responds) that catches the agent claiming
breakage was already present:

``` bash
if echo "$message" | grep -qiF "pre-existing"; then
    echo 'Blocked: Master is not broken. If tests are failing, the breakage is yours.' >&2
    exit 2
fi
```

This is a meta-circumvention: the agent can't fix the code, so instead it tries
to redefine the problem as not-a-problem.

### Why hooks work better than instructions

Instructions are suggestions. Hooks are physics. The agent can decide to ignore
an instruction under pressure; it cannot decide to ignore a hook that returns
exit code 2. The key properties:

-   **The action fails atomically** -- the file is never written, the command
    never runs. There is no partial state to clean up.
-   **The error message reaches the agent** -- it sees *why* the action was
    blocked and can adjust. Good error messages name the circumvention pattern
    explicitly.
-   **Ratchets allow incremental improvement** -- comparing old vs new means the
    agent doesn't get blocked by issues it didn't create. It only has to not
    make things worse.

## The meta-lesson

Haskell + AI is not "the AI writes Haskell for you." It's "Haskell tells you
where the AI went wrong, faster and louder than any other language." The type
system is a supervision tool -- it lets one human oversee AI output at scale.

Investment in types pays off *more* in the AI era, not less.

# Signoff

Hi, thanks for reading! You can reach me via email at <justin@jle.im>, or at
twitter at [\@mstk](https://twitter.com/mstk)! This post and all others are
published under the [CC-BY-NC-ND
3.0](https://creativecommons.org/licenses/by-nc-nd/3.0/) license. Corrections
and edits via pull request are welcome and encouraged at [the source
repository](https://github.com/mstksg/inCode).

If you feel inclined, or this post was particularly helpful for you, why not
consider [supporting me on Patreon](https://www.patreon.com/justinle/overview),
or a [BTC donation](bitcoin:3D7rmAYgbDnp4gp4rf22THsGt74fNucPDU)? :)

\end{document}

\end{document}
